\chapter{实验结果}

\section{训练结果}

\figfull[.96\textwidth]{baseline_training_loss.pdf}{正式模型完整训练历史。浅色为每次记录，深色为滚动中位趋势；橙色虚线是现场部署的第 19,000 步保存点。纵轴为对数坐标。}

训练损失从 0.067703 降至 0.000671，最低记录为 0.000499。曲线在前期快速下降，后期继续缓慢降低，说明训练过程真实发生，模型逐步更接近示教动作。\ev{E08} 现场部署保存点位于训练历史中段，并不是最低损失位置。

\plain{这张曲线只能说明“机器人越来越会模仿老师录下的动作”，不能说明“机器人在新桌面上一定能成功”。要证明后者，需要独立测试和逐次成功记录。}

\section{数据规模与质量结果}

正式训练数据的 25 个唯一仓库全部可读取，1,051 条轨迹和 879,852 帧的元数据一致；三路训练相机字段、14 维状态和 14 维动作在各仓库中一致。全量数值检查没有发现非有限值、全零动作、轨迹内时间戳不连续和帧号中断。\ev{E04,E05}

这些结果支持“数据管线具有较好数值完整性”，但不支持“视频全部无损坏”“任务语义全部正确”或“训练没有泄漏”。尤其是当前只有训练划分，无法检查训练与测试之间是否重复。

\section{训练条件与失败风险的对照}

\figfull[.94\textwidth]{prompt_condition_gap.pdf}{任务指令条件覆盖。无盖条件的 6 个仓库、158 条轨迹中，没有一个明确要求“无盖时跳过拧盖”。}

正式数据里确实存在无盖瓶轨迹，但无盖仓库仍全部使用无条件“拧盖”任务指令；只有 1 个仓库、19 条轨迹明确写出“存在瓶盖时才拧”。\ev{E11} 这会让模型同时看到“无盖画面”和“继续拧盖”的监督信号，与现场观察到的空拧行为一致。

\limitbox{这是一个明确的数据—指令风险，不是完整因果证明。仍需在相同瓶子、相同位置下成对测试“有盖/无盖”，并记录模型是否进入旋拧阶段。}

方向相关轨迹 404 条、翻转条件 127 条，说明团队已经针对瓶身方向进行过采集；但没有逐帧朝向标签和分条件自主测试，因而无法解释倒抓是视觉判断、抓取目标选择还是动作执行误差，也不能给出倒抓发生率。

\section{现场结果}

\begin{table}[H]
\centering
\caption{现场结果与证据等级}
\begin{tabularx}{\textwidth}{>{\bfseries}p{42mm}p{34mm}X}
\toprule
观察内容 & 当前表述 & 为什么这样表述\\
\midrule
完整取瓶、旋盖和分箱循环 & 已在现场展示中观察到 & 有观察记录和部署链路，缺完整录像\\
连续工作时间 & 超过 1 小时（估计） & 无起止时间表与中断记录\\
平均处理速度 & 约 2 瓶/分钟（估计） & 无逐瓶计数和有效工作时间定义\\
初步泛化 & 对随机桌面位置和一定条件变化显示适应迹象 & 无固定条件、多次重复和未见瓶型清单\\
稳定成功率 & 该指标尚未记录 & 不能由训练损失或估计吞吐量反推\\
\bottomrule
\end{tabularx}
\end{table}

结果支持“系统已经超出单次动作演示，具备完整任务和长程循环的现场表现”；不支持“已经达到某个百分比稳定成功率”或“所有瓶型都能泛化”。这是当前报告最重要的结论边界。

\section{注意力结果}

8,223 个注意力样本显示，模型在三路相机之间的中位份额约为顶部 23.2\%、左腕 33.2\%、右腕 43.5\%。腕部视野在接近瓶身、瓶口和夹爪的阶段更受关注，这与任务需要精细近景的直觉一致。\ev{E12}

然而，注意力样本没有成功/失败标签。选择一张热图只能展示记录能力，不能通过颜色浓淡判定机器人是否真的理解“有没有瓶盖”。

\section{强化学习离线结果}

\figfull[.91\textwidth]{rlt_offline_validation.pdf}{强化学习连续轮次的离线验证结果。蓝线为动作误差，橙线为价值判断误差；没有同条件基础模型和真机结果。}

从轮次 28 的初始记录到轮次 32 的最终记录，验证动作误差从约 0.1350 降至约 0.1259，相对下降约 6.7\%；价值判断误差在 3.1 左右波动，没有同步单调改善。\ev{E16} 这说明动作优化模型在固定离线验证上的拟合有所改善，同时价值学习仍不稳定。

\plain{强化学习部分现在像“在练习册上进步了”，但还没有参加同一套真机考试。只有把基础模型和强化学习模型放到相同瓶子、相同起点、相同次数下比较，才能说它提高了机器人能力。}

\section{本阶段最好结果}

当前能够确认的最好结果不是一个正式百分比，而是一组分级事实：

\begin{itemize}
  \item 完整取瓶—拧盖—瓶/盖分箱循环已在现场观察到；
  \item 连续运行超过 1 小时、约 2 瓶/分钟是现场估计；
  \item 正式模型训练、保存和部署链路可追溯；
  \item 训练误差显著下降，但没有验证/测试或正式真机指标；
  \item 强化学习离线动作误差改善约 6.7\%，但尚未证明真机增益。
\end{itemize}

\resultbox{本阶段最稳妥的成果表述是：项目已经建立完整数据与部署闭环，机器人在现场表现出完成双臂长程瓶子分拣任务的能力和初步泛化迹象；稳定成功率和标准吞吐量仍待正式测试。}

